\chapter{Theory}

\section{Category over an object}
Let $E$ be an object in a category $\mathcal{A}$.
$(\mathcal{A} \downarrow E)$ is the category of objects over $B$ in $\mathcal{A}$, we denote \\ an object by $(A, \alpha:A \to B)$.

A morphism between two objects $(A, \alpha)$ and $(A', \alpha')$ is a morphism $f: A \to A'$ such that $\alpha' f = \alpha$.

\section{Connected limits}

\begin{defi}
    A connected limit is a limit over a connected category. A category is connected if it is inhabited and any two objects can be connected by a zigzag of morphisms
    \[
    \begin{tikzcd}
        x_0 = y_1 \arrow[rd, leftrightarrow] &   & y_3 \arrow[ld, leftrightarrow] \arrow[rd, leftrightarrow] &        & y_n=x_1 \arrow[ld, leftrightarrow] \\
                          & y_2 &                                  & \cdots &             
        \end{tikzcd}
    \]
    where each symbol $\leftrightarrow$ is a morphism $y_i \to y_{i+1}$ or  $y_{i+1} \to y_i$.
    \end{defi}
    \begin{ex}  
        Some examples
        \begin{itemize}
            \item The pullback is a connected limit.
            \[
            % https://tikzcd.yichuanshen.de/#N4Igdg9gJgpgziAXAbVABwnAlgFyxMJZABgBoBGAXVJADcBDAGwFcYkRyQBfU9TXfIRTlSxanSat2AJm68QGbHgJERVGgxZtEIAMzdxMKAHN4RUADMAThAC2SMiBwQk0jZO0gscyzfuIRJxdENwktdgArAy4gA
    \begin{tikzcd}
        & 2 \arrow[d, "j"] \\
    1 \arrow[r, "i"] & 3               
    \end{tikzcd}
            \]
            \item The equalizer is a connected limit.
            \[
            % https://tikzcd.yichuanshen.de/#N4Igdg9gJgpgziAXAbVABwnAlgFyxMJZABgBpiBdUkANwEMAbAVxiRAEYQBfU9TXfIRTtyVWoxZsATNzEwoAc3hFQAMwBOEALZIyIHBCQjxzVohBYQ1OAAssqnEgC07Hms07Eeg0et2HviAMdABGMAwACvx4BGzqWAo2jtT0pmwAVrJcQA
    \begin{tikzcd}
        1 \arrow[r, "i", shift left] \arrow[r, "j"', shift right] & 2
        \end{tikzcd}\]
        \end{itemize}
    \end{ex}
    
    \begin{rem}\label{connected limits computed on the base}
    The connected limits in the category $(\mathcal{A} \downarrow E)$ can be computed in $\mathcal{A}$. In fact we can provide to the objects in $\mathcal{A}$ a morphism to form an object in $(\mathcal{A} \downarrow E)$. In the zigzag below we chose for example $(\varphi_1)(q_1)$, but we could use whatever morphism of the type $(\varphi_i)(q_i)$. The connection guarantees that all the morphisms $q_1 \ldots q_n$ involved in the limit are also morphisms in $(\mathcal{A} \downarrow E)$, ie $(\varphi_i)(q_i) = (\varphi_1)(q_1) \ \forall i \in {1,\ldots,n} $.
    % https://tikzcd.yichuanshen.de/#N4Igdg9gJgpgziAXAbVABwnAlgFyxMJZABgBoAmAXVJADcBDAGwFcYkQAxACgEYBKEAF9S6TLnyEUPUgGZqdJq3bdyA4aOx4CRchXkMWbRJy6qhIkBk0SiZACz7FRkAFFzG8dpR29NA0uNuMDULK09JZBlZR0N2AB04gGMoCBwEdUsxLQiyYhiAkAAZIXkYKABzeCJQADMAJwgAWyQyEBwIJGkFWMCuLAEaRnoAIxhGAAUsm2MsMGxYdxB6ppaadqQokFGwKCQAWjsATj8neLiGOrQACywAfR4QQZGxyesvEFn5tgzl5sQu9aITbbXaIPYyVpDUYTKbvT5YBYnHogBIXa53ciLX5IXRtDr-J7Q17hdjwxHdArcABWIVqDT+uMBwJgOyQR0JL1hkg+cwRbCRBVR9EuN1uMix9Jxa3xAFYBc5uABrAYgKGct7csnfCzYxByvFsjkwjWk3nk-wKriMWlLSWIHwGoE0EFs46q57GkkzM38inOIUiu6EH52gBs0tWWxZoPBkI9xOypq+jz97AAjvcJStEOHHV0XUDWhb07dMUaE9MecmQ9nc4DcWrPYnvcn5SXxc7oxtiDW-nX8Q6CxC28YM4Ry1yk3ys32I06o6ywaGi6djFwA+j7nwuBn+CnGLNnFB6HArmV9-HJ8Y6lhylccCVBEA
    \[
        \begin{tikzcd}
            L \arrow[dd, "q_1" description, bend right] 
              \arrow[rddd, "q_2" description, bend left] 
              \arrow[rrdd, "q_3" description, bend left] 
              \arrow[rrrrdd, "q_n" description, bend left] 
              \arrow[dddd, "(\varphi_1)(q_1)"', dashed, bend right=65] 
              & & & & \\
            & & & & \\
            F(1) 
              \arrow[leftrightarrow, rd, "F(1\leftrightarrow2)" description] 
              \arrow[dd, "\varphi_1" description, bend right=49] 
              & 
              & F(3) 
                \arrow[leftrightarrow, ld, "F(2\leftrightarrow3)" description] 
                \arrow[lldd, "\varphi_3" description, bend left=49] 
                \arrow[leftrightarrow, rd, "F(2\leftrightarrow ..)" description] 
              & 
              & F(n) 
                \arrow[lllldd, "\varphi_n" description, bend left=51] \\
            & F(2) 
              \arrow[ld, "\varphi_2" description, bend right] 
              & 
              & \cdots 
                \arrow[leftrightarrow, ru, "F(.. \leftrightarrow n)" description] 
              & \\
            E & & & & 
            \end{tikzcd}
    \]
    \end{rem}

    \begin{rem}
        In particular one can affirm that, for each object $E$ in $\mathcal{A}$, $(\mathcal{A} \downarrow E)$ has all the connected limits that $\mathcal{A}$ has.
    \end{rem}
    
    \begin{ex}
        The product is not a connected limit, in fact we can see that the following category is not connected.
        \[
            % https://tikzcd.yichuanshen.de/#N4Igdg9gJgpgziAXAbVABwnAlgFyxMJZABgBpiBdUkANwEMAbAVxiRAEYQBfU9TXfIRTtyVWoxZsATNzEwoAc3hEKXIA
    \begin{tikzcd}
        1 & 2
        \end{tikzcd}
    \]
    We try to compute the product directly in $(\mathcal{A} \downarrow E)$. Let $(D,\delta)$ and $(D',\delta')$ be two objects in $(\mathcal{A} \downarrow E)$, the product is given by the following pullback
    \[
    % https://tikzcd.yichuanshen.de/#N4Igdg9gJgpgziAXAbVABwnAlgFyxMJZARgBoAGAXVJADcBDAGwFcYkQBBAHS7wFt4AfQBCHAOQgAvqXSZc+QigDMFanSat24qTJAZseAkXKliahizaJOO2QYVEATKfMarIAKK29cw4uTOjq6W7F6SajBQAObwRKAAZgBOEHxIJiA4EEiO0gnJqYhkGVmIKuoh1jxMaAAW9BI0jPQARjCMAAq+DtaJWFE1ON5JKWk0mUjEuSDDBc7FSGUWmpVc1XVD+dljJQAsNEvuVYy19CCNLW2d9kY9fQMbI4XbSHvlyyBHJxJTMwvPiK9GFgwO4oBBmM1GGwaDUYPQoOxICCzhl6FhGIiCNCQE1Wh0ujccTB4oNwpIgA
    \begin{tikzcd}
        & D\times_ED' \arrow[ld] \arrow[rr] &                                  & D' \arrow[ld, "\delta'"'] \arrow[ldd, "\delta'"] \\
    D \arrow[rr, "\delta"] \arrow[rrd, "\delta"'] &                                   & E \arrow[d, equal] &                                                  \\
        &                                   & E                                &                                                 
    \end{tikzcd}
    \]
    which doesn't correspond to the product $D\times D'$ in $\mathcal{A}$.
    
    
    
    \end{ex}

\section{The pullback functor and its adjunction}

Given a category $\mathcal{A}$ with pullbacks and a morphism $p: E \to A$ in $\mathcal{A}$, one can consider functors pullback and shrink
\[
\begin{tikzcd}
    (\mathcal{A} \downarrow B)
  \arrow[bend left=40]{r}[description]{p^*}
  & (\mathcal{A} \downarrow E)
  \arrow[bend left=40]{l}[description]{p_!}
\end{tikzcd}
\]
defined as follows:
\begin{itemize}
\item $p^*: (\mathcal{A} \downarrow E) \to (\mathcal{A} \downarrow B)$ is the pullback functor
\end{itemize}
\[
% https://tikzcd.yichuanshen.de/#N4Igdg9gJgpgziAXAbVABwnAlgFyxMJZAZgBoAGAXVJADcBDAGwFcYkQBRAHS7wFt4AfQBCAQRABfUuky58hFAFYK1Ok1btxUmdjwEiygIyqGLNohDDJ0kBl3yiZYzVMaLHazrn6U5FS-VzEC0bO28FZD9nNTN2K21bWT0Iw38YtxBPRPsfZAAmNNcgyVUYKABzeCJQADMAJwg+JD8QHAgkVPSgtAaAK0E8rPrGpDJW9sQCrvY0IYamxBa20YDYix6IfsM5kcRO5cnVjJ4mNAALeh2FgBYaA+UQRnoAIxhGAAUkhwtGGBqcEBHIInRjnS4JYYLABsdwmAHY7vQsIx2Hx6Gg4MsaHAzlh-kgEdN1gA9ABUJQkQA
\begin{tikzcd}
    A \arrow[d, "\alpha"] & {} \arrow[r, "p^*", maps to, shift right=7] & {} & E\times_BA \arrow[rr, "\mathrm{proj}_2"] \arrow[d, "\mathrm{proj}_1"] &  & A \arrow[d, "\alpha"] \\
    B                     &                                             &    & E \arrow[rr, "p"]                                   &  & B                    
    \end{tikzcd}
\]
\[
% https://tikzcd.yichuanshen.de/#N4Igdg9gJgpgziAXAbVABwnAlgFyxMJZAFgBoAGAXVJADcBDAGwFcYkQBRAHS7wFt4AfQBCAQRABfUuky58hFAFYK1Ok1btxUmdjwEiAdhU0GLNohCiA5JOkgMu+UQAcxtWfbdeWAXBHXbHTl9FAA2UgBGVVMNC2FA+1k9BWRwgCZo9XNOBIdglPI3GOytOzzkogiirM0bbUTHEOQ06o8LXKSnFABmVtiQDsaCyMy2kHiJVRgoAHN4IlAAMwAnCD4kQpAcCCRlEAAjGDAoJABabs3i9h4+ehwAC2W+YDRVgCsJQQiElbWkKq2O0QZHc-R4TDQ93oP1W60QLUBSBBVws4MYkPoNhojHoh0YAAVOiEQMssDN7jgYX9EJttv8TDVUVxbg8ni93p80iBsbiYASiQoSWSKVS4QC6fCGWNFqKkL1EZLQdkbndHs9XhAPoI0nU7L84bSgfLDsc5ZtGFgwNkoPQ4PdptyldcuDAAB5YOA4OAAAgAhN6eHA0PQAMYwAPeRiwYCLCSyxB7CXIxn2R04vGEobsRgwRaU+r6uU0CV7E0nRAXKVg5mqtkarURLEgdN8zP5dik8n5vWwpDhBVGJ0WGUF3uIVwKgCcxfoWEY7FuaE9Oxodqwed2VeyaAAegAqeP9iURS4ptEY+OD4+nsbnqF1SgSIA
\begin{tikzcd}
    A \arrow[r, "f"] \arrow[d, "\alpha"] & A' \arrow[ld, "\alpha'"] & {} \arrow[r, "p^*", maps to, shift right=5] & {} & E\times_BA \arrow[rrdd, "\mathrm{proj}_1", bend right] \arrow[r, "\mathrm{proj}_2"'] \arrow[rrrr, "\exists ! \space \tilde{f}", dashed, bend left] & A \arrow[rd, "\alpha"] \arrow[rr, "f"] &                  & A' \arrow[ld, "\alpha'"'] & E\times_BA' \arrow[l, "\mathrm{proj}_2'"] \arrow[lldd, "\mathrm{proj}_1'"', bend left] \\
    B                                    &                          &                                             &    &                                                                                                                                                    &                                        & B                &                           &                                                                                        \\
                                         &                          &                                             &    &                                                                                                                                                    &                                        & E \arrow[u, "p"] &                           &                                                                                       
    \end{tikzcd}
\]
\begin{itemize}
    \item $p_!: (\mathcal{A} \downarrow B) \to (\mathcal{A} \downarrow E)$ is the shrink functor

\[
% https://tikzcd.yichuanshen.de/#N4Igdg9gJgpgziAXAbVABwnAlgFyxMJZABgBpiBdUkANwEMAbAVxiRABEQBfU9TXfIRRkAjFVqMWbAKLdeIDNjwEiI8uPrNWiEHL5LBRAEzrqmqTr0L+yocgDMpiVrace+gSpSOxZydpBZd2sDL2QAFlJfZwsQACFucRgoAHN4IlAAMwAnCABbJDIQHAgkNRiAgB1K2AYcOisc-KQTYtLER2K6LAY2PLo0OBKQajgACyxMnCQAVj8XHTQAfQBCRtyCxEi2pAA2edi0atr69eat6mHEOYq2Y5g6huoGOgAjB4AFG0MdbKwUsbTYJNTY3K77W6LRJcIA
\begin{tikzcd}
    D \arrow[d, "\delta"] & {} \arrow[r, "p_!", maps to, shift right=5] & {} & D \arrow[rd, "p\delta"] \arrow[d, "\delta"'] &   \\
    E                     &                                             &    & E \arrow[r, "p"]                             & B
    \end{tikzcd}
\]
\[
% https://tikzcd.yichuanshen.de/#N4Igdg9gJgpgziAXAbVABwnAlgFyxMJZABgBpiBdUkANwEMAbAVxiRABEQBfU9TXfIRQAmclVqMWbdgHJuvEBmx4CRMgEZx9Zq0QgAovL7LBRAMxjq2qXqOL+KocgAsliTrZ2lA1SgCsbta6HF4OpigA7IGSwbKhJr7IAZpWMWyGPMY+TlEp7jYgAELc4jBQAObwRKAAZgBOEAC2SGQgOBBIovnBADo9sAw4dCDUDHQARjAMAAphviB1WOUAFjh29U0t1O1I6qkeelAjIGOTM3NCJzA1a5kgG82Ie20diF1BbH0DQ3KjE1OzBKXBjXW4KB5ICwvJCuNp0LAMNiNOhoOA7ahwZZYG5IALdNhoAD6AEJ1g1HnidogAGz7ApHO4QxBRaGIAAcdOCaGOpwBFzYixWYNq5Nx21eHPxejQXymQx5-3OQIFS1Wx0x2LWiGIjNFNPFSElH2lssGdDkXAoXCAA
\begin{tikzcd}
    D \arrow[d, "\delta"'] \arrow[rr, "d"] &  & D' \arrow[lld, "\delta'"] & {} \arrow[r, "p_!", maps to, shift right=5] & {} & D \arrow[rr, "d"] \arrow[rrd, "p\delta"'] &  & D' \arrow[d, "p\delta'"] \\
    E                                      &  &                           &                                             &    & E \arrow[rr, "p"']                        &  & B                       
    \end{tikzcd}
\]
\end{itemize}

\begin{prop}\label{p_! preserves connected limits}
Let $p : E \to B$ be a morphism in $\mathcal{A}$, then $p_!$ preserves connected limits.
\begin{proof}
    By \Cref{connected limits computed on the base} connected limits are computed in $\mathcal{A}$ and $p_!$ doesn't affect objects and morphism in $\mathcal{A}$.
\end{proof}
\end{prop}
\begin{prop}
    Let $p : E \to B$ be a morphism in $\mathcal{A}$, then
    \[
\begin{tikzcd}
    (\mathcal{A} \downarrow B)
  \arrow[bend left=40]{r}[description]{p^*}
  & (\mathcal{A} \downarrow E)
  \arrow[bend left=40]{l}[description]{p_!}
  \arrow[phantom, shift left=1.5, shift right=1.5]{l}[description]{\top}
\end{tikzcd}
\]
is an adjunction.

In particular the unit is given by

\begin{minipage}{0.45\textwidth}
    \[
        \eta: \mathrm{Id}_{(\mathcal{A} \downarrow E)} \to p^*p_!
    \]
    \[
        \eta_{(D,\delta)} = \langle \delta,1_D \rangle : (D,\delta) \to (E\times_BD, \mathrm{proj}_{1})
    \]
    \end{minipage}
    \begin{minipage}{0.5\textwidth}
    \[
    % https://tikzcd.yichuanshen.de/#N4Igdg9gJgpgziAXAbVABwnAlgFyxMJZABgBpiBdUkANwEMAbAVxiRABEQBfU9TXfIRQBGUsKq1GLNgFEAOnLwBbeAH0AQpx59seAkQDMYifWatEHbrxAZdgoqIBMJqeZAyrOgfpRHn1U2kLdW4JGCgAc3giUAAzACcIJSQjEBwIJAAWANc2NE8QBKSkUTSMxEdtQsTkxFL0lJyzPMSAK1VhAqLax2oGxGzJZos0BVgGHDoumqQyMqRekAAjGDAoFLnAt2FVLWtu2b7y0oYsMDcoOjgAC3CQalu6dcQwJgYGProsBjZIc-uhkEQAoYAAPLBwHBwAAEAEJoQo4Gg6ABjGDQgA8YxgEzoYl2AD5psVEHN+qkVmskABaAxzBh0FYMAAK-D0QhA8SwEWuOABWzY2NxoS4QA
    \begin{tikzcd}
        D \arrow[rrrd, "1_D", bend left] \arrow[rd, "{\exists ! \space \langle\delta,1_D\rangle}", dashed] \arrow[rdd, "\delta"', bend right] &                                           &  &                        \\
                                                                                                                                  & E\times_BD \arrow[rr] \arrow[d] &  & D \arrow[d, "p\delta"] \\
                                                                                                                                  & E \arrow[rr, "p"]                         &  & B                     
        \end{tikzcd}
        \]
    \end{minipage}
and the counit is given by

\begin{minipage}{0.45\textwidth}
    \[
        \epsilon: p_!p^* \to \mathrm{Id}_{(\mathcal{A} \downarrow B)}
    \]
    \[
        \epsilon_{(A,\alpha)} = \mathrm{proj}_2 : (E\times_BA, \ p (\mathrm{proj}_1)) \to (A,\alpha)
    \]
    \end{minipage}
    \begin{minipage}{0.5\textwidth}
    \[
    % https://tikzcd.yichuanshen.de/#N4Igdg9gJgpgziAXAbVABwnAlgFyxMJZABgBpiBdUkANwEMAbAVxiRAFEAdTvAW3gD6AIQCCIAL6l0mXPkIoATOSq1GLNmMnTseAkSUBGFfWatEIIRKkgMOuUTJHqJ9efYSVMKAHN4RUABmAE4QvEhkIDgQSAbOamY2IQBWAgpWgSFhiADM1FFISqqmbGjpIMGh4XnROXHF5mjJAgZlFVmxkTWFLgncjGgAFnQe4kA
    \begin{tikzcd}
        E\times_BA \arrow[rr, "\mathrm{proj}_2"] \arrow[d, "\mathrm{proj}_1"] &  & A \arrow[d, "\alpha"] \\
        E \arrow[rr, "p"]                                   &  & B.                    
        \end{tikzcd}
    \]
    \end{minipage}
\begin{proof}
Let us show that this is actually an adjunction. First we want to prove that the unit is a natural transformation. Let $g : (D,\delta)  \to (D',\delta')$ be a morphism in $(\mathcal{A} \downarrow E)$, we want to prove that the following diagram commutes:
\[
% https://tikzcd.yichuanshen.de/#N4Igdg9gJgpgziAXAbVABwnAlgFyxMJZABgBpiBdUkANwEMAbAVxiRABEQBfU9TXfIRRkAjFVqMWbdgHJuvEBmx4CRAEzlx9Zq0QgAogB1DeALbwA+gCFOPPssHrSY6tql6jJrObjXZ3cRgoAHN4IlAAMwAnCFMkMhAcCCQNCR02YxgcOgtgAAp2UmNYBmyASi55SJi4xBFqJKQAZldJXRBM7NyCmSLDEuyZCqqQaNj4huS61vS9YJGx2tTGxBa090UAPQAqNAsAQjzgsoBeYzwGWGBgyq4KLiA
\begin{tikzcd}
    D \arrow[rr, "{\eta_{(D,\delta)}}"] \arrow[d, "g"] &  & E\times_BD \arrow[d, "p^*p_!(g)=\tilde{g}"] \\
    D' \arrow[rr, "{\eta_{(D',\delta')}}"]             &  & E\times_BD'                                
    \end{tikzcd}
\]

In order to do that we need to expand the action of the functor $p^*p_!$ on the morphism $g$. We have:

\[
% https://tikzcd.yichuanshen.de/#N4Igdg9gJgpgziAXAbVABwnAlgFyxMJZABgBpiBdUkANwEMAbAVxiRABEQBfU9TXfIRQBGclVqMWbdgHJuvEBmx4CRAExjq9Zq0Qh5fZYKIBmTRJ1sDi-iqElSw8dql6AotaUDVKACzmXXQ5PW2MUAFYAySDZEKMfZEinLWi2ACFucRgoAHN4IlAAMwAnCABbJDIQHAgkUQtXEBzrEvKkDWraxDNquiwGNjK6NDgakGo4AAssQpwkSIagtAB9AEIW0orEKrHEf0W2AB1D2AYcOg22xHrd-cCjk5gzujkeIs356l2AdhTLPTQx1O50uWwAbF8ur8DgCgU9zq8FK0tgtdhCYU1MlwgA
\begin{tikzcd}
    D \arrow[r, "g"] \arrow[d, "\delta"] & D' \arrow[ld, "\delta'"] & {} \arrow[r, "p_!", maps to, shift right=5] & {} & D \arrow[rd, "p\delta"] \arrow[r, "g"] & D' \arrow[d, "p\delta'"] \\
    E                                    &                          &                                             &    &                                        & B                       
    \end{tikzcd}
\]

\[
% https://tikzcd.yichuanshen.de/#N4Igdg9gJgpgziAXAbVABwnAlgFyxMJZAJgBoAGAXVJADcBDAGwFcYkQBRAHS7wFt4AfQBCAERABfUuky58hFAGYK1Ok1btxUmdjwEiAVhU0GLNohCiA5JOkgMu+UQBsxtWfbdeWAXBHXbHTl9FAAWUgBGVVMNC2FA+1k9BWRw4mj1c04Eh2CU8jcYrJykpxQIwsz2SVUYKABzeCJQADMAJwg+JAKQHAgkIxAAIxgwKCQAWkUeovYePnocAAs2vmA0DoArCUEIhPbOpAre-sRw91j7HlhGHHp9jq7EMhOkc9mLNGuYW-obGkY9BGjAACqUQiA2lh6kscA9Dogen0jiYqhZ5osVmsNhBtoJiPCnsdkc9UR4LPVCUhlK9SRcshjlqt1lsdsR-iBAcCwY4IVCYXDtCADk8kacaSMxtSeowsGAslB6HAlnUQGTLjwYAAPLBwHBwAAEAEIDTw4Gh6ABjGCm7yMWDAeoSKmIQYk95o+xqzlAn48vLsRgwFqCuwi6k0EmDSXjRDTdUMrgLJnY1m7Dlcv3ghSQ6Gwl2uWkAdkj9CwjHYCzQev6NGVWBDAwT7DQAD0AFQu4mnHox6XN9FcGB3QTAAAUolI31+AEpnULw3SSRLRrGpjNPZqR+PrFOuDc7lY597M6Ds+x+fmJJQJEA
\begin{tikzcd}
    {} \arrow[r, "p^*", maps to, shift right=5] & {} & E\times_BD \arrow[rrdd, "\mathrm{proj}_1", bend right] \arrow[r, "\mathrm{proj}_2"] \arrow[rrrr, "\exists ! \space \tilde{g}", dashed, bend left] & D \arrow[rd, "p\delta"] \arrow[rr, "g"] \arrow[l, "{\eta_{(D,\delta)}}", bend left] &                  & D' \arrow[ld, "p\delta'"'] \arrow[r, "{\eta_{(D',\delta')}}"', bend right] & E\times_BD' \arrow[l, "\mathrm{proj}_2'"'] \arrow[lldd, "\mathrm{proj}_1'"', bend left] \\
                                                &    &                                                                                                                                                   &                                                                                     & B                &                                                                            &                                                                                         \\
                                                &    &                                                                                                                                                   &                                                                                     & E \arrow[u, "p"] &                                                                            &                                                                                        
    \end{tikzcd}
\]

Considering the diagram
\[
% https://tikzcd.yichuanshen.de/#N4Igdg9gJgpgziAXAbVABwnAlgFyxMJZABgBpiBdUkANwEMAbAVxiRABEQBfU9TXfIRRkAjFVqMWbdgHJuvEBmx4CRAEzlx9Zq0QgAogB1DeALbwA+gCFOPPssHrSY6tql6jJrObjXZ8+wFVFAAWZy1JXQ45O0V+FSFkDTUInTZ9ALiHYOQwlNdItitucRgoAHN4IlAAMwAnCFMkMhAcCCQNCTS9YxgcOgtgAAp2UmNYBn6ASi5M+sakEWo2pABmAu6QXv7BkZkxwwn+mRm5hqbEFpXEJa73EHKzhcRO6-W7qLQAPQAqNAsAIRDcpTAC8xjwDFgwHKs1i8wu72uYQ+bGMpjoOAAFnVTMA0A0AFZcCxqGIKBFrZbtRAAVg293RmJxeIJEGJFhE5Nq5yQ9NaNIAbAzPk8LijrsLUXo0OMYJM6NyQJTLtSqSAGFgwFEoHQ4FiyiUuEA
\begin{tikzcd}
    D \arrow[rr, "{\eta_{(D,\delta)}}"] \arrow[d, "g"] \arrow[rrd, dashed] &  & E\times_BD \arrow[d, "p^*p_!(g)=\tilde{g}"]                              &  &                          \\
    D' \arrow[rr, "{\eta_{(D',\delta')}}"]                                 &  & E\times_BD' \arrow[rr, "\mathrm{proj}_2'"] \arrow[d, "\mathrm{proj}_1'"] &  & D' \arrow[d, "p\delta'"] \\
                                                                           &  & E \arrow[rr, "p"]                                                        &  & B                       
    \end{tikzcd}
\]
we can prove the naturality in the left diagram by the universal property of the pullback in the right. Since $(\mathrm{proj}_1$, $\mathrm{proj}_2)$ are jointly monic, by the following compositions
\begin{itemize}
    \item $(\mathrm{proj}_2')(\tilde{g})(\eta_{(D,\delta)}) = (g)(\mathrm{proj}_2)(\eta_{(D,\delta)}) = (g)(1_D) = g$
    \item $(\mathrm{proj}_2')(\eta_{(D',\delta')})(g) = (1_{D'})(g) = g$
    \item $(\mathrm{proj}_1')(\tilde{g})(\eta_{(D,\delta)}) = (\mathrm{proj}_1)(\eta_{(D,\delta)}) = \delta$
    \item $(\mathrm{proj}_1')(\eta_{(D',\delta')})(g) = (\delta')(f) = \delta$
\end{itemize}
we conclude that ($\tilde{g}) (\eta_{(D,\delta)}) = (\eta_{(D',\delta')}) (g)$. So the naturality condition of $\eta$ is satisfied.

Now we want to prove that the counit is a natural transfromation. Consider $f : (A,\alpha) \to (A',\alpha')$ a morphism in $(\mathcal{A} \downarrow B)$, we want to prove that the following diagram commutes:
\[
% https://tikzcd.yichuanshen.de/#N4Igdg9gJgpgziAXAbVABwnAlgFyxMJZABgBpiBdUkANwEMAbAVxiRAFEAdTvAW3gD6AIQCCIAL6l0mXPkIoyARiq1GLNlx5Z+cYSIDkEqSAzY8BIgCZyK+s1aIQYydLNyrpZdTvrHBiSowUADm8ESgAGYAThC8SNYgOBBIAMzeag4gEUaRMXGIZInJiIrp9mxoAgCEaAB6AFQAFBEAlAC83HgMsMAR4iDUDHQARjAMAAoy5vIgUVjBABY4OVl5SIVJ8WW+INwwaNgMBALAjSKk3IxoC3Qt-S6rsUilRanbmXsHWEdgJ2f6F04Vxu+juAXEQA
\begin{tikzcd}
    E\times_BA \arrow[d, "p_!p^*(f)=\tilde{f}"'] \arrow[rr, "{\epsilon_{(A,\alpha)}}"] &  & A \arrow[d, "f"] \\
    E\times_BA' \arrow[rr, "{\epsilon_{(A',\alpha')}}"]                                &  & A'              
    \end{tikzcd}
\]

As similarly done above, we need to expand the action of the functor $p_!p^*$ on the morphism $f$, we have:

\[
% https://tikzcd.yichuanshen.de/#N4Igdg9gJgpgziAXAbVABwnAlgFyxMJZAFgBoAGAXVJADcBDAGwFcYkQBRAHS7wFt4AfQBCAQRABfUuky58hFAFYK1Ok1btxUmdjwEiAdhU0GLNohCiA5JOkgMu+UQAcxtWfbdeWAXBHXbHTl9FAA2UgBGVVMNC2FA+1k9BWRwgCZo9XNOBIdglPI3GOytOzzkogiirM0bbUTHEOQ06o8LXKSnFABmVtiQDsaCyMy2kHiJVRgoAHN4IlAAMwAnCD4kQpAcCCRlEAAjGDAoJABabs3i9h4+ehwAC2W+YDRVgCsJQQiElbWkKq2O0QZHc-R4TDQ93oP1W60QLUBSBBVws4MYkPoNhojHoh0YAAVOiEQMssDN7jgYX9EJttv8TDVUVxbg8ni93p80iBsbiYASiQoSWSKVS4QC6fCGWNFqKkL1EZLQdkbndHs9XhAPoI0nU7L84bSgfLDsc5ZtGFgwNkoPQ4PdptyldcuDAAB5YOA4OAAAgAhN6eHA0PQAMYwAPeRiwYCLCSyxB7CXIxn2R04vGEobsRgwRaU+r6uU0CV7E0nRAXKVg5mqtkarURLEgdN8zP5dik8n5vWwpDhBVGJ0WGUF3uIVwKgCcxfoWEY7FuaE9Oxodqwed2VeyaAAegAqeP9iURS4ptEY+OD4+nsbnqF1SgSIA
\begin{tikzcd}
    A \arrow[r, "f"] \arrow[d, "\alpha"] & A' \arrow[ld, "\alpha'"] & {} \arrow[r, "p^*", maps to, shift right=5] & {} & E\times_BA \arrow[rrdd, "\mathrm{proj}_1", bend right] \arrow[r, "\mathrm{proj}_2"'] \arrow[rrrr, "\exists ! \space \tilde{f}", dashed, bend left] & A \arrow[rd, "\alpha"] \arrow[rr, "f"] &                  & A' \arrow[ld, "\alpha'"'] & E\times_BA' \arrow[l, "\mathrm{proj}_2'"] \arrow[lldd, "\mathrm{proj}_1'"', bend left] \\
    B                                    &                          &                                             &    &                                                                                                                                                    &                                        & B                &                           &                                                                                        \\
                                         &                          &                                             &    &                                                                                                                                                    &                                        & E \arrow[u, "p"] &                           &                                                                                       
    \end{tikzcd}
\]
\[
% https://tikzcd.yichuanshen.de/#N4Igdg9gJgpgziAXAbVABwnAlgFyxMJZAJgBoAGAXVJADcBDAGwFcYkQBRAHS7wFt4AfQBCAQRABfUuky58hFAGYK1Ok1btxUmdjwEiAVhU0GLNohCiA5JOkgMu+UQBsxtWfbdeWAXBHXbHTl9FAAWUgBGVVMNC2FA+1k9BWRyNxjzEASHYJSI9PVMyVUYKABzeCJQADMAJwg+JHyQHAgkcPdYkB4mNAALegS6hqQyFrbEDoz2HsZ++hsaRnoAIxhGAAUkpwtarDK+nCH6xsQ08aaTQpmuPnocPtq+YDR6gCsJQWIQJdX1rccIRAewOR20IGGp2arVGVw8FmqxxGiGUF0QY2mFh4dweTxe70+33BkKQ5xhKJoazAUCQinOjCwYEyUHocD6pR+nUyPBgAA8sHAcHAAAQAQmFPDgaHoAGMYBLvIxYMBqhIkacjGjXC16FhGOw7mhBW0aGysNUjohNZj7IJRerSTRyR0qTTEABaOlwrpoAAU2Puj2erwgH0EEQAlJzlmtNtsgSDDg6KWiXTBqbTzja-QHccGCeGrFGJJQJEA
\begin{tikzcd}
    {} \arrow[r, "p_!", maps to, shift right=5] & {} & E\times_BA \arrow[r, "\mathrm{proj}_2"'] \arrow[rrrr, "\exists ! \space \tilde{f}", dashed, bend left] \arrow[rrd, "p(\mathrm{proj}_1)"', bend right] & A \arrow[rd, "\alpha"] \arrow[rr, "f"] &   & A' \arrow[ld, "\alpha'"'] & E\times_BA' \arrow[l, "\mathrm{proj}_2"] \arrow[lld, "p(\mathrm{proj}_1')", bend left] \\
                                                &    &                                                                                                                                                       &                                        & B &                           &                                                                                       
    \end{tikzcd}
\]

Since $\epsilon_{(A,\alpha)}$ is the $\mathrm{proj}_{2,E\times_BA}$, the composition $(f)(\epsilon_{(A,\alpha)})=(\epsilon_{(A',\alpha')})(\tilde{f})$ holds by the definition of $\tilde{f}$. So also the naturality of $\epsilon$ is satisfied.

Now we can check that the unit and counit fulfil the first triangle identities. Let $(A,\alpha)$ be an object in $(\mathcal{A} \downarrow B)$, we have:
\[ ((p^*\epsilon) \circ (\eta p^*))_{(A,\alpha)} = (p^*\epsilon)_{(A,\alpha)}  (\eta p^*)_{(A,\alpha)} = p^*(\mathrm{proj}_{2,E\times_BA})  (\eta_{(E\times_BA,\mathrm{proj}_1)} ) = \circledast
\]

Considering first how $p^*$ acts on the morphism $\mathrm{proj}_{1,E\times_BA}$
\[
% https://tikzcd.yichuanshen.de/#N4Igdg9gJgpgziAXAbVABwnAlgFyxMJZABgBpiBdUkANwEMAbAVxiRAFEAdTvAW3gD6AIQCCIAL6l0mXPkIoATOSq1GLNmMnTseAkQCMpfSvrNWiEEIlSQGHXKIBmZdVPqL17bL0oALC9UzNgkVGCgAc3giUAAzACcIXiQyEBwIJCVA9xBuRjQACzoAAgAKNASAKwFgJS4eLH44YRFxAEoQagY6ACMYBgAFGV15EDiscPycTxB4xKRDVPTETLdzHM48wunZpMQUtPnXNTXyiCrgQzq+QVFxbYTd50WkfxA4fKwYqcQAViOgixoAB6ACoOqk6FgGGxeHQ0HADuIKOIgA
\begin{tikzcd}
    E\times_BA \arrow[rd, "{\alpha (proj_{2,E\times_BA})}"'] \arrow[rr, "{proj_{1,E\times_BA}}"] &   & A \arrow[ld, "\alpha"] & {} \arrow[r, "p^*", maps to, shift right=5] & {} \\
                                                                                                 & B &                        &                                             &   
    \end{tikzcd}
\]
\[
% https://tikzcd.yichuanshen.de/#N4Igdg9gJgpgziAXAbVABwnAlgFyxMJZABgBpiBdUkANwEMAbAVxiRAFEAdTvAW3gD6AIQAUXHln5xhAQQCUIAL6l0mXPkIoAzOSq1GLNuL6ChMpSpAZseAkQCsu6vWatEIc8tU2NRAOxO+q5G3CbSZhbe6nYoACykAEx6LobuQpFWaraayPFayQZuHEp6MFAA5vBEoABmAE4QvEhkIDgQSACMzoVsaA0AVgLACaTGkqZioePSMkJyiiDUDHQARjAMAApZvu51WOUAFjgZ9Y3N1G1IOkGpIFMMsMB9EIPDo1NSsooLS6vrWz4YiAGDAasdqGswFArsQvCBTk1EC1LohHDcis9Xl0xp9RDjBLN5otgX9NtsgSCwcTIdDEABaLSwywIzoXdqIeLotgibiMNAHOhyESYoYjfHhGSKBS-NZkwGaEB7Q7HOEsxBdVrskZc9wit7ir7E5aygHRBVKo4nBqI7UozkpIq8hj8uhG0mm7JsSkq5nWq5spDah29Aai94SXGSt0m8kK71Ws6Ia4otHG-6xtgW8EgGkw7rBXWh4DYj6mKOqv2ogMc-O3NDR9Pyr2glUURRAA
\begin{tikzcd}
    E\times_B(E\times_BA) \arrow[rrr, "{\mathrm{proj}_{2,E\times_B(E\times_AB)}}"'] \arrow[rrrrrrr, "{\tilde{\mathrm{proj}_{2,E\times_BA}}}", bend left] \arrow[rrrrddd, "{\mathrm{proj}_{1,E\times_B(E\times_AB)}}", bend right] &  &  & E\times_BA \arrow[rdd, "{(\alpha)(\mathrm{proj}_{2,E\times_BA})}"'] \arrow[rr, "{\mathrm{proj}_{2,E\times_BA}}"'] &                  & A \arrow[ldd, "\alpha"] &  & E\times_BA \arrow[ll, "{\mathrm{proj}_{2,E\times_BA}}"] \arrow[lllddd, "{\mathrm{proj}_{1,E\times_BA}}"', bend left] \\
                                                                                                                                                                                                       &  &  &                                                                                                 &                  &                         &  &                                                                                                    \\
                                                                                                                                                                                                       &  &  &                                                                                                 & B                &                         &  &                                                                                                    \\
                                                                                                                                                                                                       &  &  &                                                                                                 & E \arrow[u, "p"] &                         &  &                                                                                                   
    \end{tikzcd}
\]

\[ \circledast = (\tilde{\mathrm{proj}_2}) (\eta_{(E\times_BA,\mathrm{proj}_1)}) \] 
and then considering the definition of $\eta$, we obtain
\[
    % https://tikzcd.yichuanshen.de/#N4Igdg9gJgpgziAXAbVABwnAlgFyxMJZABgBpiBdUkANwEMAbAVxiRAFEAdTvAW3gD6AIQCCIAL6l0mXPkIoAjKQVVajFmy48s-OMIAUWvoNEBKCVJAZseAkQDMy1fWatEHbsb2iL0m3KIlACZndTcOXysZW3lkRxDqFw13IQlVGCgAc3giUAAzACcIXiRHEBwIJAAWRLC2NBBqBjoAIxgGAAVogPcCrEyACxxIwuKkJXLKxCDJfKKSxAmK0trXeqKAKwFgJSMdE0NPfe8RU3ER+aQg6mXEGrU19zR9NE3t3aPdYRFxc1mQUYLMiTK7UNpgKClYFJcIKbZ7L6ic7-QFIYG3CYMLBgcJQOhwAYZRogQl0SGIMBMBgMG50LAMNiQHHEmFsbgwAAeWDgODgAAIAIR87hwNB0ADGMD5AB5XhAtjtSAiTD9lPDPirxAA+C5jRDoqZlcHkgC09mBzTanW6dl6-SGLLqTzeiuVJ2RFHEQA
\begin{tikzcd}
    E\times_BA \arrow[rrrd, "1_{E\times_BA}", bend left] \arrow[rd, "{\exists ! \space \langle\mathrm{proj}_{1,E\times_BA},1_{E\times_BA}\rangle}", dashed] \arrow[rdd, "{\mathrm{proj}_{1,E\times_BA}}"', bend right] &                                                                                &  &                                                  \\
                                                                                                                                                                                         & E\times_B(E\times_BA) \arrow[rr] \arrow[d] &  & E\times_BA \arrow[d, "{p(\mathrm{proj}_{1,E\times_BA})}"] \\
                                                                                                                                                                                         & E \arrow[rr, "p"']                                                             &  & B                                               
    \end{tikzcd}
\]

\[ \circledast = (\tilde{\mathrm{proj}_2}) (\langle \mathrm{proj}_{1,E\times_BA}, 1_{E\times_BA} \rangle) \] 

In order to conclude we need to check some compositions so that, by the universal property of the pullback $E\times_BA$ in the right of the following diagram, we obtain the identity morphism.
\[
% https://tikzcd.yichuanshen.de/#N4Igdg9gJgpgziAXAbVABwnAlgFyxMJZABgBpiBdUkANwEMAbAVxiRAFEAdTvAW3gD6AIQCCIAL6l0mXPkIoAzOSq1GLNlx5Z+cYQApNfQaICUEqSAzY8BIgBZl1es1aIO3I7tHnp1uUQA2R1UXNjFJX1lbFAB2YOd1N0NtY3CLKyj5ZABWUgBGFQTXECEfSxkbLNyAJkK1YvYJFRgoAHN4IlAAMwAnCF4kMhAcCCQ8p3q2AB40PoArAWBx5J1hEUk8xZXU8QA+EGoGOgAjGAYABQr-Nx6sVoALHDLe-sHqEaQgkMTLecXljwpLzrA4gI6nC5XaIgW4PJ7UU5gKBIAC0CmIERALwGiCGH0Q1QmoTcm2A22B4lB4LOlz80IYMC68JAiORiHRmOxY3eowJhxONKh8hhd0ezz6OPGw15Dm+xUBDFgwFmEAWwEJ5LW4kp-IhtMybAZTNBrKQHIsXMQUvxX2pkLpwthYoRMCRqPRRJ+KrVAK0qyEBkB-pEJkpnIlSEJ0qQuTlbD03EYaHudBMem9iw1QZ2Zl1godbCdT3Drz50fZnuKGfVpE1oh1YIF9oNN1FxYtEYr5djRTYiYYyboVKb+sqhsZ7e6ndl+KUjb1QvHxpLOJnvNtI8XrbhJtdbI9cbc1d9ni14tLX3xPcmR+HC4LbiNxYo4iAA
\begin{tikzcd}
    E\times_BA \arrow[rrr, "{<\mathrm{proj}_{1,E\times_BA},1_{E\times_BA}>}"'] \arrow[rrrrrdd, "{\mathrm{proj}_{1,E\times_BA}}"', bend right] \arrow[rrrr, "1_{E\times_BA}", bend left] &  &  & E\times_B(E\times_BA) \arrow[r] \arrow[rrrr, "{\tilde{\mathrm{proj}_2}}", bend left] \arrow[rrdd, "{\mathrm{proj}_{1,E\times_B(E\times_BA)}}"', bend right] & E\times_BA \arrow[rd, "{(\alpha)(\mathrm{proj}_{2,E\times_BA})}"'] \arrow[rr, "{\mathrm{proj}_{2,E\times_BA}}"'] &                  & A \arrow[ld, "\alpha"] & E\times_BA \arrow[l] \arrow[lldd, "{\mathrm{proj}_{1,E\times_BA}}"', bend left] \\
                                                                                                                                                                      &  &  &                                                                                                                                                        &                                                                                                & B                &                        &                                                                        \\
                                                                                                                                                                      &  &  &                                                                                                                                                        &                                                                                                & E \arrow[u, "p"] &                        &                                                                       
    \end{tikzcd}
\]

\[ \circledast = (\tilde{\mathrm{proj}_2}) (\langle \mathrm{proj}_{1,E\times_BA}, 1_{E\times_BA} \rangle) = \mathrm{Id_{(E\times_BA,\mathrm{proj}_1)}}\] 

We want to check that composition of the projections of $E\times_BA$ with $\circledast$ is the same of composing those projection with $\mathrm{Id_{(E\times_BA,\mathrm{proj}_1)}}$.  Following the first composition
\[ 
     (\mathrm{proj}_{2,E \times_BA})(\tilde{\mathrm{proj}_2}) (\langle \mathrm{proj}_1, 1_{E\times_BA} \rangle) = (\mathrm{proj}_{2,E \times_BA}) (\mathrm{proj}_{2,E\times_B(E\times_BA)}) (\langle \mathrm{proj}_1, 1_{E\times_BA} \rangle) =
\]
\[
    = (\mathrm{proj}_{2,E \times_BA}) (1_{E\times_BA}) = \mathrm{proj}_{2,E \times_BA} = (\mathrm{proj}_{2,E \times_BA}) (\mathrm{Id_{(E\times_BA,\mathrm{proj}_1)}})
\]
and the second composition
\[
    (\mathrm{proj}_{1,E \times_BA})(\tilde{\mathrm{proj}_2}) (\langle \mathrm{proj}_1, 1_{E\times_BA} \rangle) = (\mathrm{proj}_{1,E\times_B(E\times_BA)})(\langle \mathrm{proj}_1, 1_{E\times_BA} \rangle) =   
\]
\[   
    \mathrm{proj}_{1,E \times_BA} = (\mathrm{proj}_{1,E \times_BA}) (\mathrm{Id_{(E\times_BA,\mathrm{proj}_1)}})
\]

Finally we can check the second triangle identities. Let $(D,\delta)$ be an object in $(\mathcal{A} \downarrow E)$, we have:
\[
((\epsilon p_!) \circ (p_! \eta))_{(D,\delta)} = ((\epsilon p_!)_{(D,\delta)} (p_! \eta))_{(D,\delta)} = (\epsilon_{p_!(D,\delta)}) p_! (\langle \delta, 1_D \rangle) =
\]
\[ 
    = (\epsilon_{(D,p \delta)}) p_! (\langle \delta, 1_D \rangle) = \mathrm{proj}_{2,E\times_BD} (\langle \delta, 1_D \rangle) = 1_D = \mathrm{Id}_{(D,\delta)}
\]
\[
% https://tikzcd.yichuanshen.de/#N4Igdg9gJgpgziAXAbVABwnAlgFyxMJZABgBpiBdUkANwEMAbAVxiRAFEAdTvAW3gD6AIQAiIAL6l0mXPkIoATOSq1GLNmMnTseAkSUBGFfWatEIIRKkgMOuUTJHqJ9efYSVMKAHN4RUABmAE4QvEhkIDgQSAbOamYg3Lx0OAAWQbzAaCEAVuICwEpcPFj8cMIi4laBIWGIAMzUUUhKqqZsaNUgwaHhTdENce3m2RA5BbHFfIKiVVrdtTH9LUOuNtywDDh0HuJAA
\begin{tikzcd}
    E\times_BD \arrow[rr, "{\mathrm{proj}_{2,E\times_BD}}"] \arrow[d, "{proj_{1,E\times_BD}}"] &  & D \arrow[d, "p\delta"] \\
    E \arrow[rr, "p"]                                                                          &  & B                     
    \end{tikzcd}
\]

So we can conclude that the adjunction holds.
\end{proof}
\end{prop}
\begin{rem}
$p^*$ is a right adjoint functor, so it preserves limits. $p_!$ is a left adjoint functor, so it preserves colimits.
\end{rem}

We need the following definition in order to claim some properties.

\begin{defi}
    Let $\gamma$ be a natural transformation between two functors $F,G : \mathcal{C} \to \mathcal{D}$, we say that $\gamma$ is a cartesian natural transformation if for every morphism $f : A \to B$ in $\mathcal{C}$, the following diagram is a pullback:
\[
% https://tikzcd.yichuanshen.de/#N4Igdg9gJgpgziAXAbVABwnAlgFyxMJZABgBpiBdUkANwEMAbAVxiRADEAKAQQEoQAvqXSZc+QigBM5KrUYs2AcR78hI7HgJFpARln1mrRCGUAhVcJAYN4omT3UDC413ODZMKAHN4RUADMAJwgAWyQyEBwIJB1HeSMQAB1ErzoQkLoAfWBuAUFLINDw6iikAGY4wzYufwsA4LDECsjoxGk5KuNk1PSs03z6osRYlqR2pwTlWvcBIA
\begin{tikzcd}
    F(A) \arrow[rr, "\gamma_{A}"] \arrow[d, "F(f)"] &  & G(A) \arrow[d, "G(f)"] \\
    F(B) \arrow[rr, "\gamma_B"]                     &  & G(B)                  
    \end{tikzcd}
\]
ie the naturality diagrams are pullback.
\end{defi}

\begin{propr}
The unit $\eta$ and the counit $\epsilon$ are cartesian.
\begin{proof}

First we prove that $\eta$ is cartesian. Let $g : (D,\delta) \to (D',\delta')$ be a morphism in $(\mathcal{A} \downarrow E)$. Considering the commutative diagram below, since both the bottom and the total diagrams are pullback, it follows that the top diagram is a pullback.
\[
% https://tikzcd.yichuanshen.de/#N4Igdg9gJgpgziAXAbVABwnAlgFyxMJZABgBpiBdUkANwEMAbAVxiRAFEAdTvAW3gD6AIQAiIAL6l0mXPkIoATOSq1GLNmMnTseAkTIBGFfWatEHbn0GiA5BKkgMOuUSVHqJ9eZF2tjmbryJKQKxmpmHPbasnqKIWGmbEISKjBQAObwRKAAZgBOELxIZCA4EEgGHuFs3Lx0OAAWebzAaAUAVuICClEg+YVISqXliADMVYnmtfVNLW0Qnd2+Dv1FiCVlgxNejgB6AFRoAgCEABTpAJQAvJZYDLDA6eIg1Ax0AEYwDAAKAS7mDBgORwvVWFWomzG2wi6ReIDenx+f1iIDyWHSDRBfjBiCGkIALNCapw6o1mq0Ol0DMtcgU1uNhkgAKxE8xobiwBg4Oh2V4fL6-ZwotEYrErOlIQmMxAs1STRygiWISrS2WfMBQJAANhKngi7M4nO5ioG6whIyl6s1iAAtDrWSBpmS5pSBAY4QiBcj5Kj0ZiUuIgA
\begin{tikzcd}
    E\times_BD \arrow[rr, "\mathrm{proj}_2"] \arrow[d, "p^*p_!(g)=\tilde{g}"] \arrow[dd, "\mathrm{proj}_1"', bend right=60] &  & D \arrow[d, "g"'] \arrow[dd, "p\delta", bend left=60] \\
    E\times_BD' \arrow[rr, "\mathrm{proj}_2'"] \arrow[d, "\mathrm{proj}_1'"]                                                &  & D' \arrow[d, "p\delta'"']                             \\
    E \arrow[rr, "p"]                                                                                                       &  & B                                                    
    \end{tikzcd}
\]

Considering the following commutative diagram, since the right diagram is a pullback, as we just proved, and the total diagram is a pullback, we can conclude that the left diagram is a pullback, as stated.
\[
% https://tikzcd.yichuanshen.de/#N4Igdg9gJgpgziAXAbVABwnAlgFyxMJZAJgBoAGAXVJADcBDAGwFcYkQBRAHS7wFt4AfQBCAERABfUuky58hFABYK1Ok1btxUmdjwEiZAIyqGLNok49+QsQHJJ0kBl3yiy4zVMaLo+9qeyegrI5KQeamaafo7OcvoooVSe6uYgWqowUADm8ESgAGYAThB8SKEgOBBIhsmRFjx89DgAFoV8wGjFAFYSgsQOBcWliGQVVYgAzLXeIA1Nre2dED190YMlZTSVSKNeqWgAegBUaIIAhAAUWQCUALxWWIywwFkSIDSM9ABGMIwACoFXBZGDB8jgBiAihtEDUxkgphEZll3iBPj9-oD4iBClgss1wf4ocNlHCRtNUjwYDh6IJgBdfKQeLBGNTbNc3oShkgAKxbcYkvbsZGc6G80nlQX1LhUml00SMrjM6nsiFEpAk7aTGg-MBQJAAWgmEpS7EMtN8bw+31+AJcWJxeIJjjViDFmthOr1k2NdRAZvEVvRtriClRoIJlAkQA
\begin{tikzcd}
    D \arrow[d, "g"] \arrow[rr, "{\eta_{(D,\delta)}}"] \arrow[rrrr, "1_D", bend left] &  & E\times_BD \arrow[rr, "\mathrm{proj}_2"] \arrow[d, "p^*p_!(g)=\tilde{g}"] &  & D \arrow[d, "g"'] \\
    D' \arrow[rr, "{\eta_{(D',\delta')}}"] \arrow[rrrr, "1_{D'}"', bend right]        &  & E\times_BD' \arrow[rr, "\mathrm{proj}_2'"]                                &  & D'               
    \end{tikzcd}
\]

Now we wan to prove that $\epsilon$ is cartesian. Let $f : (A,\alpha) \to (A',\alpha')$ be a morphism in $(\mathcal{A} \downarrow B)$. Considering the following commutative digram, since the bottom diagram is a pullback and the total diagram is a pullback, we can conclude that the top diagram is a pullback, as stated.
\[
% https://tikzcd.yichuanshen.de/#N4Igdg9gJgpgziAXAbVABwnAlgFyxMJZABgBpiBdUkANwEMAbAVxiRAFEAdTvAW3gD6AIQCCIAL6l0mXPkIoATOSq1GLNmMnTseAkTIBGFfWatEHbn0GiA5BKkgMOuUSVHqJ9eZF2tjmbryJKQKxmpmHPbasnqKIWGmbEISKjBQAObwRKAAZgBOELxIZCA4EEgGHuFs3DBo2AwEAsAAFCKk3IxoABZ0AJTiALzcvHQ43Xm8wGgFAFbiAgpRIPmFSEql5YgAzFWJ5rX1WI1gzW02HZxdvTYDw5yj45PTcwsKvg6rRYglZet7XkcAgAhGgAHoAKhaOT69zwDFgwBy4hA1AYdAARjAGAAFAIucwMGA5HDLL4Vah-HYAiI5VEgdFY3H42IgPJYdLdUl+cmIDZUgAsNJqDzGEymMwg8wEBg+uQK312myQAFZhQcrgwenQ7GjMdi8c5WezOdzPgqkELlYg1ap9o4yRbEJVrbasWAoEgAGwlTwRTpa3qOtY-SlbK3uz2IAC0PvVIBGYuekulBnpjINLPkbI5XJS4iAA
\begin{tikzcd}
    E\times_BA \arrow[rr, "{\epsilon_{(A,\alpha)}=\mathrm{proj}_2}"] \arrow[d, "p_!p^*(f)=\tilde{f}"] \arrow[dd, "\mathrm{proj}_1"', bend right=60] &  & A \arrow[d, "f"'] \arrow[dd, "\alpha", bend left=60] \\
    E\times_BA' \arrow[rr, "{\epsilon_{(A',\alpha')}=\mathrm{proj}_2'}"] \arrow[d, "\mathrm{proj}_1'"]                                              &  & A' \arrow[d, "\alpha'"']                             \\
    E \arrow[rr, "p"]                                                                                                                               &  & B                                                   
    \end{tikzcd}
\]
\end{proof}
\end{propr}

\section{The monad associated to the adjunction}

Consider the monad associated to the adjunction we just studied, we want to study some properties that will be useful.



\begin{propr}
The monad $(p^*p_!, \eta, p^*\epsilon p_!)$ is cartesian.
\begin{proof}
We need to prove that the category $(\mathcal{A} \downarrow E)$ has pullback, the monad preserves pullbacks, and both the multiplication $p^* \epsilon p_!$ and the unit $\eta$ are cartesian.

The category $(\mathcal{A} \downarrow E)$ has pullback, since pullbacks are connnected limit and hence are computed in $\mathcal{A}$, which has pullbacks. The monad $p^*p_!$ preserves pullbacks, in fact $p^*$ is a right adjoint so it preserves limits, therefore preserves pullbacks, and $p_!$ preserves pullbacks because (cita all'inizio). We already proved that the unit $\eta$ is cartesian, so it remains to prove that the multiplication $p^* \epsilon p_!$ is cartesian. Since $\epsilon$ is cartesian and $p^*$ preserves pullbacks as shown above, we can conclude that the multiplication is cartesian.
\end{proof}
\end{propr}

\begin{defi}
A monad $(T,\eta,\mu)$ on a category $\mathcal{X}$ with finite coproducts is said to be essentially nullary if the morphism $[T(!_X),\eta_X]$ is a strong epimorphism for each object $X$ in $\mathcal{X}$, where $!_X$ is the unique morphism from the initial object of $\mathcal{X}$ to $X$.
\end{defi}

Now we claim a lemma which will be useful to prove some theorems.
\begin{lemma}
    A category $\mathcal{X}$ is protomodular id and only if for each pullback
    \[
    % https://tikzcd.yichuanshen.de/#N4Igdg9gJgpgziAXAbVABwnAlgFyxMJZABgBpiBdUkANwEMAbAVxiRACEByAHW7wFt4AfXYBBEAF9S6TLnyEUAJnJVajFm3FSZ2PASJkAjKvrNWiDp0nSQGXfKLLj1UxovtJqmFADm8IqAAZgBOEPxIZCA4EEiGLurmIADW1kGh4YjKUTGIAMzxZmy8AEYwOHSpICFhsdTRSPkgcAAWWIE4tWqFFmgg1Ax0pQwACrJ6CiDBWD7NHdpV6Q11OXFNre2dA0Oj9voWUzMdBW5NldUZWfWIkS1tHYirWzAjYw7707Nni9fLSFm3Gwe-UGzx2cj2kw+cwoEiAA
\begin{tikzcd}
    B'\times_BA \arrow[rr, "k"] \arrow[d, shift right] &  & A \arrow[d, "p"', shift right] \\
    B' \arrow[rr, "\beta"] \arrow[u, shift right]      &  & B \arrow[u, "s"', shift right]
    \end{tikzcd}
    \]
    the pair $(k,s)$ is stronly epimorphic.
\end{lemma}
The proof is out of the scope of this thesis, but we can find it in (citazione) %\cite{Borceux}.

\begin{theo}
If $\mathcal{X}$ is a pointed protomodular category with finite coproducts and $(T,\eta,\mu)$ is a monad on $\mathcal{X}$ with cartesian $\eta$, then the monad is essentially nullary.
\begin{proof}
    For each object $X$ in $\mathcal{X}$ we can consider the following diagram:
    \[
    % https://tikzcd.yichuanshen.de/#N4Igdg9gJgpgziAXAbVABwnAlgFyxMJZABgBpiBdUkANwEMAbAVxiRAA0QBfU9TXfIRQAmclVqMWbACoAKdgEpuvEBmx4CRMgEZx9Zq0QhiyvusFFRu6vqlG5xJV3EwoAc3hFQAMwBOEAFskMhAcCCRtG0lDEAAdWJgcOgB9Th4ffyDEUVDwxABmKIM2eMSUk3SQP0CI6jCkQpA4AAssbxxaiWL7WQBCAD1FEGoGOgAjGAYABX4NIRBfLDdmjsrqrMb6xEim1vbO2xi5XtSlEfHJmfNNI0Xl1ZV1pBytkJa2ju2iuxATznOJtNZhZbksVqYqplgnU8jl3vsvl0fgN-iBRoCrgIbgswasKFwgA
\begin{tikzcd}
    X \arrow[rr, "\eta_X"] \arrow[d, "!^X"', shift right] &  & T(X) \arrow[d, "T(!^X)"', shift right] \\
    0 \arrow[rr, "\eta_0"] \arrow[u, "!_X"', shift right] &  & T(0) \arrow[u, "T(!_X)"', shift right]
    \end{tikzcd}
\]
Since $\mathcal{X}$ is protomodular then by the lemma above, the pair $(\eta_X,T(!_X))$ is a strongly epimorphic. It is a standard fact that given $f$ and $g$ morphism, if $(f,g)$ is a strongly epimorphic pair, then $[f,g]$ is a strong epi. So we conclude that $[\eta_X,T(!_X)]$ is a strong epi, ie the moand is essentially nullary. (citazione) %\cite{https://arxiv.org/pdf/1409.4219) ALGEBRAICALLY COHERENT CATEGORIES
\end{proof}
\end{theo}

% https://tikzcd.yichuanshen.de/#N4Igdg9gJgpgziAXAbVABwnAlgFyxMJZAFgBoBmAXVJADcBDAGwFcYkQAdDgW3pwAsATt2ABhCDACOAXwAUaAPoBGUooBMAShDTS6TLnyEUaitTpNW7AILbdIDNjwEiABlM0GLNok48+QkQBRSVlJDWlbPUdDV1I1M09LHzQAPQAqWS5eAWFgYLkwrR0og2djOISLb3t02SsNAF4rLjxueAUAIUDI+30nIxIKjyr2VIys-1zxKTlFFXUNIrsHUoHyUhdKr3ZRAHIelf6iFU3h7Z9RbTMYKABzeCJQADNBCG4kFRAcCCQ3c3OQJIei83kgTF8fohPnB+FgnjgkABaT6JarqYGvd6IcHfD40GFwhFQs5JezKEA0Rj0ABGMEYAAU+jEfIIsLd+AjiiAQVj1hCwVyeUgyPyoYLMUgAKw0XGIFzi0GIPmykUE+FIlEjZK1BYNFpYRiwYDqCKUml0xnRMogRgwdUK3kyyGq2Hq4n-UljeTKRr6w0wY3KU0280MpnW1nszl2IWIEWy6UetHpYChcJ6jh4f3AGQYxXxyEANgdUqdSELlKwYGqUHoMJuFKT7C4MAAHlg4Dg4ABCAAE0eeEsQAHYy3H8a6iZqAYIlI2qbSw1ajCBIxy81jR6KXYSNSTqoIFGoN0gt7K+ajmxw2x2u72+wgzYvLat2GuEZXq+xa-WoFdpEAA
\begin{tikzcd}
    & C \arrow[rdd, "r1"', shift right] \arrow[rdd, "r_2", shift left] \arrow[ldd, "\exists ! s"', dashed] &                                                                          & C' &                                                                           \\
    &                                                                                                      &                                                                          &    &                                                                           \\
p^*(\mathrm{Eq}(q)) \arrow[d] \arrow[rr, "p^*(p_2)=\tilde{p_2}", shift left] \arrow[rr, "p^*(p_1)=\tilde{p_1}"', shift right] &                                                                                                      & p^*(A)=A\times_BE \arrow[d] \arrow[rr, "p^*{(q)}=\tilde{q}"] \arrow[ruu] &    & {p^*(\mathrm{Coeq}(p_1,p_2))} \arrow[d] \arrow[luu, "\exists! t", dashed] \\
\mathrm{Eq(q)} \arrow[rr, "p_2", shift left] \arrow[rr, "p_1"', shift right]                                                  &                                                                                                      & A \arrow[rr, "q"]                                                        &    & {\mathrm{Coeq}(p_1,p_2)}                                                 
\end{tikzcd}